\section{RANDOMIZED WEIGHTED BH AND THE SHARPER BOUND}
\label{app:wbhcounter}

Theorem~\ref{thm:design} gives WCS the bound $\alpha m_0/m$. The example below shows that BH applied to the same randomized weighted $p$-values of Equation~\eqref{eq:randp} need not satisfy that bound under the same design. It does not show a failure of FDR control at the nominal level $\alpha$, which remains open.

\paragraph{A three-node example.}
Take two normal nodes $a$ and $b$ and one anomaly $c$, with scores $s_c<s_a<s_b$, acquisition probabilities $\rho(a)=\rho(c)=0.3$ and $\rho(b)=1$, and weights $w=1/\rho$. One of $a,b$ is chosen uniformly as the null test node, the anomaly is always tested ($m=2$, $m_0=1$), the other normal is acquired with its probability and enters the reference, and $\alpha=0.5$, so the BH thresholds are $0.25$ and $0.5$. If $a$ is tested, the reference is $\{b\}$ and $p_a$ and $p_c$ are independent and uniform on $[3/13,1]$; conditional on $a$ being tested, the expected FDP is $0.35^2/2+0.025\times0.65=0.0775$, since FDP is $1/2$ when both are rejected and $1$ when $a$ alone is. If $b$ is tested and $a$ was acquired (probability $0.3$), then $p_b\le3/13<0.25$ and $p_c\ge0.5$, so $b$ alone is rejected and FDP $=1$. If $b$ is tested and $a$ was not acquired, the reference is empty, $p_b$ and $p_c$ are independent uniforms, and the conditional expected FDP is $0.25$. Hence
\[
\mathrm{FDR}=\tfrac12(0.0775)+\tfrac12\{0.3+0.7(0.25)\}=\tfrac{221}{800}=0.27625>\alpha m_0/m=0.25 .
\]
For randomized WCS with homogeneous pruning on the same design, integrating exactly over the randomized $p$-values and the pruning coin gives conditional expected FDPs $0.020625$ when $a$ is tested, and $1$ and $0.1875$ when $b$ is tested with and without a reference, so its FDR is $723/3200=0.2259375$, below the bound as Theorem~\ref{thm:design} requires; a Monte Carlo estimate ($400{,}000$ draws) agrees. Both procedures are below the nominal $0.5$.

\paragraph{Scope of the search.}
The example came from a protocol fixed before execution: exact evaluation of $20{,}000$ random finite designs within the assumptions of Theorem~\ref{thm:design} (two to six normals, at most four test nodes, $\rho$ on a grid from $0.01$ to $1$, $\alpha\in\{0.05,0.1,0.2,0.3,0.5\}$), followed by local search. Exact values were cross-checked against an independent implementation and Monte Carlo. Randomized weighted BH exceeded $\alpha m_0/m$ in $618$ of the $20{,}000$ randomly generated designs (a descriptive count that depends on how designs were generated, not a failure rate), at every $\alpha$ tested, by at most $0.0009$ at $\alpha\le0.1$ and $0.026$ at $\alpha=0.5$; no design exceeded the nominal $\alpha$, which leaves nominal validity unresolved rather than established, deterministic weighted BH never exceeded $\alpha m_0/m$, and every exceedance required an anomaly in the test set. A separate five-node design with the largest excess (three normals and two anomalies, excess $0.027$), not the three-node example above, was replicated $2$ to $25$ times with the same probabilities and score blocks, and under that replication scheme the excess was no longer detectable (Monte Carlo, $20{,}000$ draws per size). This does not rule out counterexamples in larger populations. The example separates the sharper guarantees of the two procedures; it is not evidence about the benchmark graphs, where no clear violation was detected (Section~\ref{sec:results}).
